\section{Native discovery when important borrower variables are hidden}
\label{exec14:cxx-latent}

A credit officer may observe an offer, recorded arrears and eventual recovery
while missing the household shock that affected all three. A useful discovery
component must preserve that possibility. Returning a convenient directed
graph can change the subsequent intervention question even when every number
and variable name survives the software boundary. The relevant native choices
are therefore compared by their graph meanings, input assumptions and failure
behavior, as well as by language and deployment requirements.

Three inspected families provide concrete alternatives. aGrUM 3.2.0 has a
native FCI path returning a partial ancestral graph (PAG). rCausalMGM 1.0.1
exposes C++ FCI-Stable through an R interface for mixed data. MIIC 2.0.4 also
has R and C++ sources, but carries information-based orientation scores and
consensus fields that need their own interpretation. The legacy MIIC repository
explicitly points to the newer R package. aGrUM's implementations named Miic
and SimpleMiic are separate source paths; the shared method name does not make
these packages interchangeable
\citep{exec14_native_agrum,exec14_native_rcmgm,exec14_native_miic}.

These are selected primary-source findings at revisions dated no later than
7 October 2026, acquired later. The repository searches and known alternatives
are bounded. A language-filtered search returned no FCI hits despite the
inspected native implementation, demonstrating why a zero search count cannot
establish absence. Other returned MIIC-like names concerned games, coursework
and unrelated applications. The broader scholarly search still has
uninspected leads. Thus the comparison supports a prototype choice without
claiming a complete census or executed engine qualification.

\subsection{aGrUM: preserve the native PAG before asking an effect question}

The inspected learner requires the FCI algorithm for \texttt{learnPAG()} and
rejects missing database values in that route. Its FCI implementation performs
skeleton discovery, collider orientation, possible-D-Sep pruning, endpoint
reset and repeated orientation rules. It also exposes conditioning-set and
path restrictions. Those settings belong in the result: shortening a search
is a scientific change as well as a resource choice. Source comments and
upstream test fixtures do not establish completeness or finite-sample coverage
of this particular implementation \citep{exec14_native_agrum}.

The native PAG distinguishes a circle, a tail and an arrowhead at each
endpoint. Its accessor \texttt{markAt(src,dst)} refers to the mark at the
\emph{destination}. An adapter must bind a mark to its node identifier, not
assume that a method's first argument names the endpoint being returned. A
circle records unresolved orientation within the represented class; a tail
or arrowhead asserts more. The statistical assumptions supporting that class
still need review, including the distinction between latent confounding and
selection. The earlier FCI/RFCI discussion supplies the theoretical context
\citep{exec13_latentr_colombo2012}.

The generic conversion is consequential. \texttt{PAG::toMixedGraph()} checks
only whether each endpoint is an arrowhead. Consequently $X\mathbin{\circ\!\to}Y$
and $X\to Y$ both become the same directed arc; circle--circle and tail--tail
both become an undirected edge. Two arrowheads become two opposite arcs.
The subsequent FCI mixed-structure bridge calls a routine that removes
opposing arcs, sometimes removing both in its fallback. This route is useful
for a different graph representation, but it cannot serve as a lossless PAG
export \citep{exec14_native_agrum}.

An exact representation calculation enumerates the nine local pairs from
the three endpoint marks taken in pairs. The conversion maps them to
four generic outputs with class sizes $4,2,2,1$. The point is not that all nine
pairs form globally valid PAGs. It is that the representation itself cannot
recover the original marks after conversion. The proposed boundary therefore
serializes native endpoints directly and retains the raw native graph as
evidence. No conversion test can by itself verify a discovered graph's causal
meaning \citep{exec14_native_oracle}.

\subsection{rCausalMGM: the R contract remains part of the native method}

The documented \texttt{fciStable} entrypoint accepts numeric continuous
variables and factor or character categorical variables. It drops rows with
missing values, can start from a supplied undirected skeleton, and defaults
to majority collider orientation and possible-D-Sep search. It also offers
max-p, conservative and separating-set rules. These alternatives can return
different statements; they are not merely display settings. A skeleton
screen that removes a true adjacency cannot be justified by the speed of the
later FCI pass \citep{exec14_native_rcmgm}.

The inspected C++ code initializes circles, applies an RFCI-style R0 pruning
step, optionally searches possible-D-Sep and delegates final orientations.
Turning off possible-D-Sep changes the route. Neither the method name nor a
successful call would establish equivalence to another package's RFCI output.
The exported R graph list includes edge strings, graph type, algorithm and
ambiguous triples. The native edge type has tail, arrowhead and circle fields,
plus star and null states. An adapter should handle supported states explicitly
and refuse the others, preserving ambiguous triples and configuration.

The inspected string parser splits an edge into exactly three space-delimited
components and recognizes selected endpoint spellings. Stable internal IDs
are therefore preferable to embedding business labels with spaces in edge
strings. A tiny round-trip fixture must check both directions, partial and
bidirected edges and unsupported tokens before admission. The source review
does not claim that this fixture has been run. Likewise, silently dropping
accounts changes the observation law: either justify the complete-case target
and selection premises, or refuse that requested analysis.

\subsection{MIIC: retain scores, modes and unresolved orientation conventions}

The successor MIIC package declares version 2.0.4, Rcpp linkage and imports
including ppcor, scales and stats. Its README requests a C++14 compiler and
the inspected build fragment delegates OpenMP flags to R. This supports an
R-hosted native candidate, not a claim of a standalone C++ service or successful
local build \citep{exec14_native_miic}.

The default latent setting is \texttt{orientation}, distinct from
\texttt{yes} and \texttt{no}. The output contains adjacency codes, two
orientation-score columns, information quantities, an optional shuffle-based
confidence ratio, and, under consistency iteration, graph-cycle and consensus
fields. A directed edge with an uncertain tail is different from the package's
stronger causal label. That label is conditioned on allowing latent variables
and disabling propagation in the inspected parser. These fields must travel
together; an adjacency matrix alone cannot preserve their meaning.

There is a concrete conversion question to settle before an adapter is enabled.
The manual describes a probability-matrix entry as referring to its row
endpoint, while the summary calls the entry at $[x,y]$ \texttt{p\_x2y}.
The native code stores an $x$-to-$y$ score in its $(x,y)$ edge and flattens
the array. The final R reshaping step and an asymmetric two-node fixture are
needed to resolve the exported convention. This is a source-level ambiguity
to investigate, not a demonstrated runtime defect. The proposal retains the
raw matrices and named summary columns until that issue is settled.

An orientation score, permutation ratio or frequency among iterated graphs
is not automatically a calibrated probability that a bank intervention has
a stated effect. The linked methodological originals remain to be inspected
before adopting their theoretical guarantees. Temporal mode additionally
declares stationarity; its presence does not qualify arbitrary account panels,
policy feedback or changing bank conditions. A separate temporal contract is
required.

\subsection{The banking consequence and the proposed reuse decision}

Consider two deliberately stipulated worlds with independent accounts.
Let $U\sim\operatorname{Bernoulli}(1/2)$ and
$V\sim\operatorname{Bernoulli}(1/4)$ be independent. In the hidden-risk world,
$A=U$ and recovery is $Y=100+20U+5V$ dollars. In the direct-effect world,
$A\sim\operatorname{Bernoulli}(1/2)$ independently of $V$ and
$Y=100+20A+5V$. Both worlds generate exactly the same observed distribution:
\[
\begin{array}{c|rrrr}
(A,Y)&(0,100)&(0,105)&(1,120)&(1,125)\\
\Pr(A,Y)&3/8&1/8&3/8&1/8.
\end{array}
\]
An intervention replaces the equation for $A$, leaving the other structural
equations fixed. Its mean effect is zero in the first world and USD~20 in the
second. At an assumed USD~8 offer cost, the incremental values are USD~$-8$
and USD~12. Both the observation law and the arithmetic reproduce exactly;
the worlds are illustrative and supply no bank effect estimate
\citep{exec14_native_oracle}.

This example explains why preserving graph ambiguity can change the decision.
More records from this same law cannot distinguish the two stipulated worlds
without additional information or restrictions. A faster native search, a
language-generated edge or a graph conversion cannot supply that information.
The next scientific question concerns an admissible experiment, a measured
common cause or another defended identifying restriction.

For the first native prototype, prefer aGrUM's explicit FCI-to-PAG boundary
with direct endpoint serialization. Retain rCausalMGM through its R interface
as a mixed-data comparator and MIIC through its own score-and-graph contract
after the unresolved source issues are settled. Porting an entire R/C++ method
would create responsibility for preprocessing, graph semantics and upgrades
without a demonstrated deployment benefit. The existing pcalg path provides
another comparator; it should not be counted twice because it has native
dependencies.

Rights and maintenance also differ. The selected aGrUM files declare
LGPL-3.0-or-later or MIT. rCausalMGM declares GPL-3. MIIC's DESCRIPTION declares
GPL at least version 2, its README lists GPL-2 or GPL-3, and its bundled full
license is GPL-3. These inspected notices need file and dependency review
before redistribution. Process separation alone is not a license conclusion.
The prototype budget must cover dependency locking, wrapper maintenance,
round-trip and failure fixtures, and comparison against independent graph
oracles. None of the three engines has been executed or admitted by this
review; the evidence supports a testable reuse recommendation.
\label{exec14:cxx-latent-end}
